\begin{frame}{Accumulated feedback}

\begin{exampleblock}{\faCodePullRequest~Pull Requests as assessment architecture}
Changes are proposed, reviewed, discussed, and integrated -- or explicitly deferred
\end{exampleblock}

\begin{columns}[T,onlytextwidth]
\begin{column}{.55\textwidth}
\begin{block}{\faRepeat~Structured accumulation}
\begin{itemize}
    \item same rubric reused
    \item prior feedback revisited
    \item unresolved issues explicitly carried forward
\end{itemize}
\end{block}
\end{column}

\begin{column}{.38\textwidth}
\begin{block}{\faToolbox~Enabled by tooling}
\begin{itemize}
    \item[\faGithub] version control
    \item[\faGitAlt] persistent work traces
\end{itemize}
\end{block}
\end{column}
\end{columns}

\end{frame}
\note{
In practice, this is implemented through \emph{pull requests}.  

Each piece of work is proposed, reviewed, discussed, and then either accepted or deferred.  

So instead of submitting assignments and moving on,  
the work remains \emph{open and revisitable}.  

If an issue is raised and not resolved, it stays visible.  

That means later work is evaluated in light of earlier feedback.  

Students cannot simply move on --  
they have to \emph{revisit and improve}.  

Because the same rubric is reused,  
I can also directly see what has changed --  
and whether earlier issues have been addressed.  

So assessment becomes more focused,  
and grading is faster because I am not re-evaluating everything from scratch.  

This creates continuity across the course,  
and is much closer to how work functions in professional settings.
}

\begin{frame}{Evolving feedback}

\begin{columns}[T,onlytextwidth]
\begin{column}{.25\textwidth}
\begin{block}{\normalsize\faChalkboardUser~Instructor}
\begin{itemize} \small
    \item method
    \item implementation
\end{itemize}
\end{block}
\end{column}
\begin{column}{.25\textwidth}
\begin{block}{\normalsize\faIndustry~Industry partner}
\begin{itemize} \small
    \item domain relevance
    \item decision context
\end{itemize}
\end{block}
\end{column}
\begin{column}{.40\textwidth}
\begin{block}{\normalsize\faUser~Student}
\begin{itemize} \small
    \item knowledge transfer through PRs
    \item defending decisions
    \item critiquing peer work
\end{itemize}
\end{block}
\end{column}
\end{columns}

\begin{columns}[onlytextwidth]
\begin{column}{.50\textwidth}
\begin{exampleblock}{\faClock~Progression over time}
\begin{enumerate}
    \item Align within teams
    \item Cross-team review 
    \item Industry validation
\end{enumerate}
\end{exampleblock}
\end{column}
\begin{column}{.40\textwidth}    
\begin{alertblock}{\faNetworkWired~Audience shift}
\centering
From internal review\\
to external scrutiny
\end{alertblock}
\end{column}

\end{columns}

\end{frame}
\note{
\begin{enumerate}
    \item \emph{Who} gives feedback:
\begin{itemize}
    \item The instructor focuses on method and implementation.  
    \item The industry partner focuses on relevance and decision context.  
    \item And importantly, students themselves are a key part of the process.  

    Through PRs, they engage in \emph{knowledge transfer} --  
    by defending their own work and critiquing others.  
\end{itemize}
    \item \emph{When} does the feedback evolve:
\begin{itemize}
    \item During the modules, the focus is on aligning within teams.  
    \item In the capstone, this expands to cross-team review, bringing in a fresh perspective.  
    \item And finally, the work is exposed to the industry partner.  
\end{itemize}

    \item \emph{What this means} is the audience widens:
    \begin{itemize}
        \item Feedback moves from internal  
        to external and more critical scrutiny.  
        \item And students are expected to respond to all of these --  
        not just receive them.
    \end{itemize}


\end{enumerate}
}